% Model section for a search-and-matching model of loan contract design
% (spread + covenant tightness), structured after Section 3 of
% Ewens, Gorbenko and Korteweg (2022, JFE), "Venture Capital Contracts".
% Compile standalone with: xelatex model_setup.tex  (or pdflatex)
\documentclass[12pt]{article}
\usepackage{amsmath,amssymb,amsthm}
\usepackage[margin=1in]{geometry}
\usepackage{natbib}
\usepackage{hyperref}
\newtheorem{proposition}{Proposition}
\newtheorem{definition}{Definition}
\newcommand{\E}{\mathbb{E}}
\newcommand{\Prob}{\mathrm{Pr}}
\newcommand{\RB}{R_B}
\newcommand{\RP}{R_P}

\begin{document}

\section{Model}
\label{sec:model}

We specify a dynamic search-and-matching model of the market for bank loans
in which a lender and a borrower who meet negotiate a loan contract consisting
of an interest-rate spread and a covenant package. The structure follows
\citet{ewens2022} (hereafter EGK): agents on both sides of the market are
heterogeneous in quality, search is random and time-consuming, the lender
makes a take-it-or-leave-it offer, and the borrower's bargaining power arises
endogenously from the option to reject and continue searching. The model is
designed to address the identification problem that arises because loan terms
are chosen jointly with the identity of the counterparty: a regression of
outcomes on contract terms is contaminated by the unobserved qualities of the
bank and of the project that drove both the match and the terms.

\subsection{Setup}
\label{sec:setup}

Time is continuous and the horizon is infinite. Two populations of agents
search for each other: banks and projects (equivalently, project owners or
borrowers). There are $m_B$ banks and $m_P$ projects actively searching at any
point in time.

\paragraph{Types.}
Each bank is characterized by a quality $i$ and each project by a quality $j$.
Qualities are drawn independently from Beta distributions,
\begin{equation}
  i \sim \mathrm{Beta}(a_i, b_i), \qquad j \sim \mathrm{Beta}(a_j, b_j),
  \label{eq:types}
\end{equation}
rescaled to the support $[0, \bar{q}]$, with cumulative distribution functions
$F_B(i)$ and $F_P(j)$. The parameters $(a_i, b_i)$ and $(a_j, b_j)$ describe
the shape of the two quality distributions. Types are fixed over time and are
observed by both parties upon meeting, but are unobserved by the
econometrician.

\paragraph{Project cash flow.}
A project requires one unit of external financing and produces a terminal
cash flow $Y$ at the end of the loan period. Cash flow is log-normally
distributed,
\begin{equation}
  \log Y = \log j + \epsilon - \tfrac{1}{2}\,\sigma(i,j)^2,
  \qquad \epsilon \sim \mathcal{N}\!\left(0, \sigma(i,j)^2\right),
  \label{eq:cashflow}
\end{equation}
so that expected cash flow equals project quality, $\E[Y] = j$. The
profitability shock $\epsilon$ is realized after the contract is signed. Its
volatility depends on the qualities of both parties,
\begin{equation}
  \sigma(i,j) = \sigma_0 + w \exp(-i) + (1-w) \exp(-j),
  \label{eq:sigma}
\end{equation}
where $\sigma_0 > 0$ is a baseline volatility and $w \in [0,1]$ governs the
relative contribution of bank quality (through screening and monitoring) and
project quality to cash-flow risk. A higher-quality bank or project lowers
the volatility of cash flow but, by construction, leaves its mean unchanged.

\paragraph{Search.}
A bank meets a project at Poisson rate $\lambda_B$, and a project meets a bank
at Poisson rate $\lambda_P$. Search is undirected: the type of the
counterparty is an independent draw from $F_P$ (for a bank) or $F_B$ (for a
project), and does not depend on the searcher's own type. Agents discount the
future at rate $r$, which we also take to be the risk-free (LIBOR) rate. Upon
meeting, the two parties either sign a contract and exit the market, or part
ways and continue searching. Matched agents are replaced by new entrants of
the same type so that the type distributions of searching agents are
stationary.\footnote{As in EGK, stationarity requires that the flow of
matches is balanced across the two sides, $\lambda_B m_B p_B = \lambda_P m_P
p_P$, where $p_B$ and $p_P$ are the average match probabilities per meeting
on each side.}

\subsection{Loan contracts and payoffs}
\label{sec:contracts}

\paragraph{Contract.}
Upon meeting, bank $i$ offers project $j$ a loan contract
\begin{equation}
  c = (s, k), \qquad s \ge 0,\; k \in [0,1],
  \label{eq:contract}
\end{equation}
where $s$ is the all-in-drawn spread over the risk-free rate and $k$ is the
covenant tightness. The loan has face value one, so the contractual repayment
at maturity is
\begin{equation}
  D(s) = 1 + r + s.
  \label{eq:repayment}
\end{equation}

\paragraph{Default.}
Two events transfer control rights to the bank. Economic default occurs when
cash flow is insufficient to repay the loan,
\begin{equation}
  P_E(i,j,s) = \Prob\!\left[Y < D(s)\right]
  = \Phi\!\left(\frac{\log D(s) - \log j + \tfrac{1}{2}\sigma(i,j)^2}{\sigma(i,j)}\right),
  \label{eq:PE}
\end{equation}
where $\Phi$ is the standard normal distribution function. Technical default
occurs when a covenant is violated. Covenant tightness is defined, as in the
empirical literature, as the ex-ante probability of violation, so that
\begin{equation}
  P_T = k .
  \label{eq:PT}
\end{equation}
Conditional on the project being solvent, technical default is independent of
the realization of $Y$. The overall probability that control rights transfer
to the bank is therefore
\begin{equation}
  P_E + (1 - P_E)\, P_T .
  \label{eq:control}
\end{equation}

\paragraph{Recovery technology.}
Covenants raise the fraction of project value the bank recovers in economic
default,
\begin{equation}
  \theta(k) = \underline{\theta} + \left(\bar{\theta} - \underline{\theta}\right)
  \left(1 - \exp(-\eta\, k)\right),
  \label{eq:theta}
\end{equation}
where $\underline{\theta}$ is the baseline recovery rate with no covenants,
$\bar{\theta}$ is the maximum attainable recovery rate, and $\eta > 0$
governs how quickly tighter covenants raise recovery. When a solvent project
violates a covenant, the bank takes control and a fraction $1 - \theta_T$ of
project value is lost; $\theta_T \in (0,1]$ is the retention rate after
technical default.

\paragraph{Payoffs.}
Given a match $(i,j)$ and a contract $c = (s,k)$, the bank's expected payoff
on the loan is
\begin{equation}
  \E[\RB(i,j,c)] = (1 - P_E)\Big[(1 - P_T)\,D
    + P_T\, \E\!\left[\min\{\theta_T Y, D\}\,\middle|\, Y > D\right]\Big]
    + P_E\, \E\!\left[\theta(k)\, Y \,\middle|\, Y < D\right],
  \label{eq:RB}
\end{equation}
and the project owner's expected payoff is
\begin{equation}
  \E[\RP(i,j,c)] = (1 - P_E)\Big[(1 - P_T)\,\E\!\left[Y - D \,\middle|\, Y > D\right]
    + P_T\, \E\!\left[\max\{\theta_T Y - D, 0\} \,\middle|\, Y > D\right]\Big].
  \label{eq:RP}
\end{equation}
The owner is protected by limited liability and receives nothing in economic
default. Summing the two payoffs,
\begin{equation}
  \E[\RB] + \E[\RP] = \E[Y]
    - (1-P_E) P_T\, \E\!\left[(1-\theta_T) Y \,\middle|\, Y > D\right]
    - P_E\, \E\!\left[(1-\theta(k)) Y \,\middle|\, Y < D\right],
  \label{eq:surplus}
\end{equation}
so total match surplus equals expected cash flow net of the deadweight losses
incurred in the two control-transfer states. The contract determines both the
size of this surplus, through $P_E(s)$, $P_T = k$, and $\theta(k)$, and its
division between the two parties. Because $Y$ is log-normal, every
conditional expectation in \eqref{eq:RB}--\eqref{eq:surplus} has a closed
form, so payoffs can be evaluated without simulation.

\subsection{Matching and the contract decision}
\label{sec:matching}

\paragraph{Outside options.}
Let $V_B(i)$ denote the value to a bank of type $i$ of continuing to search,
and $V_P(j)$ the corresponding value for a project of type $j$. A bank can
always invest in the risk-free asset, so its outside option in a meeting is
$\max\{V_B(i),\, 1 + r\}$. A project is fully externally financed; if the
owner rejects an offer he receives nothing in the current period and
continues searching, so his outside option is $V_P(j)$.

\paragraph{Acceptance.}
When bank $i$ and project $j$ meet, the pair draws a common match-specific
taste shock $\xi \sim \mathcal{N}(0, \sigma_\xi^2)$, independent of types and
of the cash-flow shock. The match is consummated under contract $c$ if both
parties' surpluses, augmented by the taste shock, are positive. The
probability that a match forms is therefore
\begin{equation}
  P(i,j,c) = \Phi\!\left(\frac{\min\Big\{\E[\RB(i,j,c)] - \max\{V_B(i), 1+r\},\;
     \E[\RP(i,j,c)] - V_P(j)\Big\}}{\sigma_\xi}\right).
  \label{eq:P}
\end{equation}
The taste shock plays the role of EGK's match-specific shock: it smooths the
acceptance decision, so that the match probability is a continuous function
of the contract and of the model parameters, and it allows the model to
generate dispersion in contract terms within a type pair.\footnote{The
shock is drawn after the contract is offered and is not contractible, so it
does not enter the contract decision in \eqref{eq:cstar} except through
$P(i,j,c)$. The parameter $\sigma_\xi$ is a scale parameter: it measures the
magnitude of the non-pecuniary component of a match relative to the
pecuniary surplus.}

\paragraph{Contract choice.}
The bank proposes the contract. Anticipating the acceptance probability
\eqref{eq:P}, bank $i$ chooses the contract that maximizes its expected value
from the meeting,
\begin{equation}
  c^*(i,j) = (s^*, k^*) = \arg\max_{(s,k)}\;
  P(i,j,c)\, \E[\RB(i,j,c)] + \big(1 - P(i,j,c)\big) \max\{V_B(i), 1 + r\}.
  \label{eq:cstar}
\end{equation}
% NOTE: the working draft (eq. 14) writes the argmax over the project owner's
% objective, P * E[R_P] + (1-P) * V_P. That objective would be solved by the
% borrower. Since the bank proposes (as in EGK, where the VC proposes), the
% objective must be the bank's. Change the text above if the intended
% proposer is the borrower.
The bank trades off a larger share of surplus, obtained through a higher
spread and tighter covenants, against a lower probability that the project
accepts and against the surplus destroyed by a higher default probability and
a higher incidence of control transfer. The project's bargaining power arises
entirely from its outside option $V_P(j)$: a project that expects to meet a
better bank soon demands more surplus today. Denote the equilibrium match
probability by $P^*(i,j) = P(i,j,c^*(i,j))$.

\subsection{Value functions and equilibrium}
\label{sec:equilibrium}

The values of search satisfy the following Bellman equations, which are the
analogue of Proposition~1 in EGK with the deterministic acceptance indicator
replaced by the match probability $P^*$.

\begin{proposition}[Value functions]
\label{prop:value}
In a stationary equilibrium the value of search for a bank of type $i$ and
for a project of type $j$ satisfy
\begin{align}
  V_B(i) &= \frac{\lambda_B}{r + \lambda_B}
  \int \Big[ P^*(i,j)\, \E[\RB(i,j,c^*(i,j))]
    + \big(1 - P^*(i,j)\big) \max\{V_B(i), 1 + r\} \Big] dF_P(j),
  \label{eq:VB} \\
  V_P(j) &= \frac{\lambda_P}{r + \lambda_P}
  \int \Big[ P^*(i,j)\, \E[\RP(i,j,c^*(i,j))]
    + \big(1 - P^*(i,j)\big) V_P(j) \Big] dF_B(i).
  \label{eq:VP}
\end{align}
\end{proposition}

The factor $\lambda/(r+\lambda)$ is the expected discount applied to the next
meeting, which arrives after an exponentially distributed waiting time with
mean $1/\lambda$. Inside the integral, with probability $P^*$ the meeting
results in a loan and the agent receives the contract payoff; with
probability $1 - P^*$ the agent returns to the market and retains the value
of search.

\begin{definition}[Equilibrium]
A stationary search equilibrium consists of value functions $V_B(\cdot)$,
$V_P(\cdot)$, a contract policy $c^*(\cdot,\cdot)$, and a matching
probability $P^*(\cdot,\cdot)$ such that (i) given $V_B$ and $V_P$, the
contract $c^*(i,j)$ solves \eqref{eq:cstar} for every pair $(i,j)$;
(ii) $P^*(i,j)$ is given by \eqref{eq:P} evaluated at $c^*(i,j)$; and
(iii) $V_B$ and $V_P$ satisfy \eqref{eq:VB}--\eqref{eq:VP}.
\end{definition}

The equilibrium is computed by value-function iteration as in EGK. Starting
from an initial guess $(V_B^0, V_P^0)$, we solve \eqref{eq:cstar} and
\eqref{eq:P} on a discrete grid of type pairs, update the value functions
using \eqref{eq:VB}--\eqref{eq:VP}, and iterate until the values converge.
The contract policy $c^*(i,j)$, the match probability $P^*(i,j)$, and the
implied default, violation and recovery outcomes are then the model objects
whose moments are matched to the data.

\paragraph{Discussion.}
Three features distinguish the model from EGK and merit comment. First, in
EGK the contract affects firm value through a reduced-form function $h(c)$;
here the contract affects value through explicit default and recovery
channels, \eqref{eq:PE}--\eqref{eq:theta}, so that the ``cost'' of a higher
spread is a higher default probability and the ``benefit'' of a tighter
covenant is a higher recovery rate, bought at the price of more frequent
control transfers. Second, the split of surplus is not a separate object:
it is implied by the debt payoff structure in \eqref{eq:RB}--\eqref{eq:RP}.
Third, acceptance is probabilistic rather than deterministic. The share of
accepted matches in which one party's pecuniary surplus is negative is an
equilibrium outcome that we report alongside the estimates.

\begin{thebibliography}{1}
\bibitem[Ewens et al.(2022)]{ewens2022}
Ewens, M., A.~S. Gorbenko, and A.~Korteweg (2022).
\newblock Venture capital contracts.
\newblock \emph{Journal of Financial Economics} 143(1), 131--158.
\end{thebibliography}

\end{document}
